\documentclass{article}
\usepackage{amsmath}
\usepackage{graphicx}
\graphicspath{ {images/} }
\usepackage[spanish]{babel}
\usepackage{bbm}
\usepackage{amsthm}
\usepackage{authblk}
\usepackage{hyperref}
\usepackage[utf8]{inputenc}
\usepackage{mathrsfs}
\usepackage{amsfonts}
\usepackage{amssymb}
\usepackage{enumerate}
\usepackage[left=3cm,right=2cm,top=2cm,bottom=2cm]{geometry}
\usepackage{epsfig}
 \renewcommand{\baselinestretch}{0.93}
 \thispagestyle{empty}
 \pagestyle{empty}
\setlength{\textheight}{53pc} \setlength{\textwidth} {36pc}
\setlength{\topmargin}{-4mm}
\usepackage{multicol}
\newcommand*{\email}[1]{%
    \normalsize\href{mailto:#1}{#1}\par
    }
\setlength{\columnsep}{1cm}
\newcommand{\N}{\mathbb{N}}
\newcommand{\calM}{\cal{M}}
\newcommand{\calP}{\cal{P}}
\newcommand{\F}{\mathbb{F}}
\newcommand{\R}{\mathbb{R}}
\newcommand{\Z}{\mathbb{Z}}
\newcommand{\Q}{\mathbb{Q}}
\newcommand{\C}{\mathbb{C}}
\newcommand{\bbH}{\mathbb{H}}


\theoremstyle{plain}
\newtheorem{thm}{Teorema}[section]
\newtheorem{prop}[thm]{Proposición}
\newtheorem{lemma}[thm]{Lema}
\newtheorem{cor}[thm]{Corolario}
\theoremstyle{definition}
\newtheorem{rec}[thm]{Recuerdo}
\newtheorem{defin}[thm]{Definición}
\newtheorem{ejem}[thm]{Ejemplo}
\newtheorem{ejer}[thm]{Ejercicio}
\newtheorem{sol}[thm]{Solución}
\newtheorem{rmk}[thm]{Observación}
\author{Marcos Morales}
\affil{Universidad Finis Terrae}
\title{Primera Semana\\}






\begin{document}


\begin{picture}(400, 75)
   \put(-40,15){\includegraphics[width=5cm]{UFT.png}}
\end{picture}


\begin{center}
\huge{Sexta Semana}
\end{center}

\begin{center}
\Large{ Marcos Morales, Camila Muñoz}
\end{center}

\begin{center}
	\large{Cálculo Vectorial}
\end{center}
\begin{center}
\setlength{\parindent}{0cm}\large{Clases centradas para capacitar a los estudiantes de ingeniería en Cálculo Vectorial. \\}
\end{center}


\vspace{1cm}

\begin{center}
	\large{Contenidos:}
\end{center}

\setlength{\parindent}{2.9cm}- Funciones vectoriales\\
\phantom{abefwfewefwfefffw} - Vector normal, vector binormal y vector tangente	\\
\phantom{abefwfewefwfefffw} - Longitud de una curva	\\
\phantom{abefwfewefwfefffw} - Curvatura	\\



\vspace{6cm}




\pagestyle{empty}


\setcounter{page}{1}
\pagestyle{plain}
\setlength{\parindent}{0cm}





\newpage

\section{Funciones Vectoriales}


Ahora lo que vamos a hacer es estudiar con mayor detenimiento las funciones vecotriales, recordemos que estas son funciones de la forma $f: \mathbb{R} \to \mathbb{R}^n$, es decir, son funciones en una variable pero cuya imagen es un vector, estas funciones en dos y tres dimensiones representan curvas parametrizadas en el plano y el espacio. En estas funciones los conceptos vistos antes toman un significado diferente y es necesario tratarlos aparte. \\


\subsection{Derivadas }

\begin{defin}Dada una función vectorial $r: \mathbb{R} \to \mathbb{R}^3$ dada por $r(t) = (r_1(t),r_2(t),r_3(t))$, podemos definir la \textbf{derivada $r'$}, como
\end{defin}


\begin{align*}
    r'(t) = (r_1'(t),r_2'(t),r_3'(t))
\end{align*}

Es decir, se deriva coordenada a coordenada, esta definición corresponde a la primera columna de la matriz jacobiana de la función vectorial. Este vector también se conoce como \textbf{Vector velocidad} y al vector $r''(t)$ se le conoce como \textbf{vector aceleración} \\

\begin{figure}[h]
    \centering
    \includegraphics[width=0.5\textwidth]{Imagenes/derivada.png}
    \caption{Intepretación gráfica de $r'(t)$}
    \label{fig:etiqueta}
\end{figure}

La derivada $r'(t_0)$ representa directamente a un vector tangente a la curva en $t_0$. \\

\begin{defin}Dada una función vectorial $r(t)$ se define el vector tangente unitario a la curva en  $t$ como
\end{defin}

\begin{align*}
    T(t) = \frac{r'(t)}{||r'(t)||}
\end{align*}

Notemos que $T(t)$ es un vector unitario, es decir $||T(t)|| = 1$. La ecuación de la recta tangente a la curva en $t=t_0$ viene dada en su ecuación vectorial por

\begin{align*}
    L :  r(t_0)+ \lambda T(t_0), \lambda \in \mathbb{R}
\end{align*}

Aunque en realidad se puede usar cualquier vector tangente, no necesariamente el tangente unitario. \\

\begin{ejem}Sea la curva dada por $r(t) = (t, \sen(t), t^2)$, tenemos que
\end{ejem}


\begin{align*}
    r'(t) = \left( 1,\cos(t), 2t\right)
\end{align*}

Luego

\begin{align*}
   || r'(t) || = \sqrt{1+\cos^2(t)+4t^2}
\end{align*}

El vector tangente  será

\begin{align*}
    T(t)= \frac{\displaystyle \left( 1,\cos(t), 2t\right) }{\sqrt{1+\cos^2(t)+4t^2}}
\end{align*}

La ecuación vectorial  de la recta tangente a $r(t)$ en el punto $t=0$ viene dada por

\begin{align*}
    L &:  r(0)+ \lambda T(0), \lambda \in \mathbb{R} \\
    L &:  (1,0,0) + \lambda \left( \frac{(1,1,0)}{\sqrt{2}}\right), \lambda \in \mathbb{R} \\
\end{align*}

Adicionalmente exiten el vector normal y binormal \\

\begin{defin}Dada una función vectorial $r(t)$ se define el \textbf{vector normal}  unitario a la curva en  $t$ como
\end{defin}

\begin{align*}
    N(t) = \frac{T'(t)}{||T'(t)||}
\end{align*}

Se define el  \textbf{vector binormal} como

\begin{align*}
    B(t) = T(t) \times N(t)
\end{align*}

\begin{figure}[h]
    \centering
    \includegraphics[width=0.5\textwidth]{Imagenes/vectores.png}
    \caption{Vectores tangente, normal y binormal de un campo vectorial $f$}
    \label{fig:etiqueta}
\end{figure}



Estos vectores nos permiten determinar las ecuaciones del plano normal, el plano rectificador y el plano osculador. \\


\newpage


\subsection{Plano Osculador, plano normal y plano rectificante}



\begin{defin}Sea $P =(x_0,y_0,z_0)\in \mathbb{R}^3$ un punto de la curva $r(t)$, con $r(t_0) = (x_0,y_0,z_0)$, el \textbf{plano osculador de $r$ en el punto $P$}, es el plano que pasa por el punto $(x_0,y_0,z_0)$ y que tiene por vector normal $B(t_0)$, su ecuaci\'on viene dada por
\end{defin}

\begin{align*}
    B(t_0) \cdot ((x,y,z)- (x_0,y_0,z_0)) = 0
\end{align*}

\begin{defin}Sea $P =(x_0,y_0,z_0)\in \mathbb{R}^3$ un punto de la curva $r(t)$, con $r(t_0) = (x_0,y_0,z_0)$, el \textbf{plano normal de $r$ en el punto $P$}, es el plano que pasa por el punto $(x_0,y_0,z_0)$ y que tiene por vector normal $T(t_0)$, su ecuaci\'on viene dada por
\end{defin}

\begin{align*}
    T(t_0) \cdot ((x,y,z)- (x_0,y_0,z_0)) = 0
\end{align*}



\begin{defin}Sea $P =(x_0,y_0,z_0)\in \mathbb{R}^3$ un punto de la curva $r(t)$, con $r(t_0) = (x_0,y_0,z_0)$, el \textbf{plano rectificante de $r$ en el punto $P$}, es el plano que pasa por el punto $(x_0,y_0,z_0)$ y que tiene por vector normal $N(t_0)$, su ecuaci\'on viene dada por
\end{defin}

\begin{align*}
    N(t_0) \cdot ((x,y,z)- (x_0,y_0,z_0)) = 0
\end{align*}


\begin{figure}[h]
    \centering
    \includegraphics[width=0.5\textwidth]{Imagenes/planoss.png}\
    \caption{Descripci\'on gr\'afica de los planos normal, rectificante y osculador}
    \label{fig:etiqueta}
\end{figure}


\begin{ejer}Sea $r(t) = (\cos(t),\sin(t),t)$, determine el vector tangente, el vector normal y el vector binormal, adem\'as determine las ecuaciones del plano osculador, el plano normal y el plano rectificante en $t=0$. \\
\end{ejer}

\begin{sol}Calculamos la derivada de la curva
\end{sol}

\[
r'(t) = (-\sin t, \cos t, 1).
\]

Evaluando en $t = 0$:
\[
r'(0) = (0, 1, 1), \quad ||r'(0)|| = \sqrt{0^2 + 1^2 + 1^2} = \sqrt{2}.
\]

Por tanto, el vector tangente  es

\[
T(0) = \frac{1}{\sqrt{2}}(0, 1, 1).
\]

Para el Vector normal  Calculamos la segunda derivada

\[
r''(t) = (-\cos t, -\sin t, 0), \quad r''(0) = (-1, 0, 0).
\]

Como $r''(0)$ es ortogonal a $r'(0)$ y tiene norma 1, es el vector normal unitario es:
\[
N(0) = (-1, 0, 0).
\]

Para el vector binormal calculamos el producto cruz
\[
B(0) = \frac{1}{\sqrt{2}}(0,1,1) \times (-1,0,0) = \frac{1}{\sqrt{2}}(0,-1,1).
\]


En $t=0 $, tenemos que $r(0) = (1, 0, 0)$.

\subsubsection*{Plano osculador}
El plano osculador tiene como vector normal $B(0) = \frac{1}{\sqrt{2}}(0, -1, 1)$, luego su ecuaci\'on es

\[
\frac{1}{\sqrt{2}}(0, -1, 1) \cdot (x - 1, y, z) = 0
\]

O bien


\begin{align*}
    z-y =0
\end{align*}

\subsubsection*{Plano normal}
El plano normal tiene como vector normal  $T(0) = \frac{1}{\sqrt{2}}(0, 1, 1)$, luego su ecuaci\'on es

\[
\frac{1}{\sqrt{2}}(0,1,1) \cdot (x - 1, y, z) = 0
\]

O bien

\begin{align*}
    y+z =0
\end{align*}

\subsubsection*{Plano rectificante}

El plano rectificante tiene como vector normal  $N(0) = (-1, 0, 0)$:

\[
(-1, 0, 0) \cdot (x - 1, y, z) = 0
\Rightarrow -(x - 1) = 0
\]

O bien

\begin{align*}
    x =1
\end{align*}

\newpage









\subsection{Longitud de una curva }

Dada una funci\'on vectorial $r(t) = (r_1(t), r_2(t), r_3(t))$ y dos puntos $r(t_0) =a$ y $r(t_1) = b$ que pertenecen a la curva, es natural preguntarse cuanto mide el largo de la curva entre $a$ y $b$, para poder responder esta pregunta usaremos las mismas t\'ecnicas vistas en c\'alculo integral, vamos a sub-dividir el intervalo $(t_0,t_1)$ en $n$ distintos intervalos $(t_{i-1},t_i)$ de igual largo $\Delta t$, generando una partici\'on de la curva


\begin{figure}[h]
    \centering
    \includegraphics[width=0.8\textwidth]{Imagenes/curvavectorial.png}\
    \caption{Partici\'on de la curva }
    \label{fig:etiqueta}
\end{figure}

Luego el largo de cada segmento viene dado por $|| r(t_{i})-r(t_{i-1})||$ y entonces creamos la suma parcial

\begin{align*}
    S = \sum_{i=0}^n || r(t_{i})-r(t_{i-1})||
\end{align*}

Ahora bien,

\begin{align*}
    || r(t_{i})-r(t_{i-1})|| &= \sqrt{(r_1(t_i)-r_1(t_{i-1}))^2+(r_2(t_i)-r_2(t_{i-1}))^2+(r_3(t_i)-r_3(t_{i-1}))^2}
\end{align*}

Ahora bien, por el teorema del valor medio, para cada $r_k$, existe $c_{k,i} \in (t_{i-1},t_i)$ tal que $r_1(t_i)-r_1(t_{i-1}) = r'(c_{k,i})\Delta t$, luego


\begin{align*}
    || r(t_{i})-r(t_{i-1})|| &= \sqrt{r_1'(c_{1,i})^2+r_2'(c_{2,i})^2+r_3'(c_{3,i})^2}\Delta t \\
    &= ||r'(c_i)|| \Delta t
\end{align*}

Donde $c_i = (c_{1,i}, c_{2,i}, c_{3,i}) $, luego la suma queda

\begin{align*}
    S = \sum_{i=0}^n || r'(c_i)||\Delta t
\end{align*}

Hciendo $\Delta t \to 0$, obtenemos la f\'ormula integral

\begin{align*}
    L = \int_{t_0}^{t_1} || r'(t)|| dt
\end{align*}

\begin{rmk}Esta demostraci\'on no es del todo rigurosa porque la forma en la que hemos llegado a la integral es un poco ambigua, en el sentido de que no deber\'ia conicidir exactamente con la integral de Riemann cl\'asica, sin embargo una demostraci\'on m\'as rigurosa tomar\'ia demasiado tiempo. \\
\end{rmk}

\textbf{Resultado:} Sea $r: \mathbb{R} \to \mathbb{R}^3$ función vectorial, la longitud de la curva entre $t=a$ y $t=b$ viene dada por

\begin{align*}
    L(a,b) = \int_a^b ||r'(t)|| dt
\end{align*}

En caso de que se especifiquen los puntos de la curva en lugar de los valores de $a$ y $b$, hay que encontrarlos utilizando la definición de la función. \\


\begin{ejem}Calcular la longitud de la curva parametrizada por
\end{ejem}
\[
\vec{r}(t) = \langle \cos(t), \sin(t), t \rangle, \quad \text{para } t \in [0, 2\pi].
\]

\begin{sol}Primero derivamos
\end{sol}
\[
r'(t) = ( -\sin(t), \cos(t), 1 )
\]

Calculamos la norma

\[
|| r'(t) || = \sqrt{(-\sin(t))^2 + (\cos(t))^2 + 1^2} = \sqrt{\sin^2(t) + \cos^2(t) + 1} = \sqrt{1 + 1} = \sqrt{2}.
\]

Por lo tanto, la longitud de la curva es:
\[
L = \int_0^{2\pi} || r'(t) || \, dt = \int_0^{2\pi} \sqrt{2} \, dt = \sqrt{2} \cdot (2\pi - 0) = 2\pi\sqrt{2}.
\]

\begin{ejer}Determinar la longitud de la curva $r(t) = (1,t,t^2)$, entre los puntos $(1,0,0)$ y $(1,1,1)$ \\
\end{ejer}

\textit{Pista: Para calcular la integral considere la sustituci\'on $t=\frac{1}{2}\tan(\theta)$, tambi\'en use que}

\begin{align*}
    \int \sec^3(x)dx = \frac{1}{2}\sec(x)\tan(x)+\frac{1}{2}\ln|\sec(x)+\tan(x)|
\end{align*}

\newpage

\subsection{Curvatura}


\begin{defin}Sea $r: \mathbb{R} \to \mathbb{R}^3$ función vecotrial, la curvatura de $r$ en el punto $t$ se define por
\end{defin}

\begin{align*}
    \tau(t) = \frac{||r'(t) \times r''(t)||}{|| r'(t)||^3}
\end{align*}






Las interpretaciones geométrica es directas del nombre. la curvatura representa que tan pronunciada es la forma en que la funci\'on vectorial se curva en $t$. \\

\begin{ejer}Calcular la curvatura de la curva
\end{ejer}

\[
r(t) = \langle t, t^2, 0 \rangle.
\]

En $t=1$ \\

\begin{sol}Tenemos que
\end{sol}

\[
\tau(t) = \frac{||r'(t) \times r''(t)||}{||r'(t)||^3}.
\]

Calculamos las derivadas primera y segunda
\[
r'(t) = \langle 1, 2t, 0 \rangle, \quad r''(t) = \langle 0, 2, 0 \rangle.
\]

El producto cruz de los vectores es
\[
r'(t) \times r''(t) =
\begin{vmatrix}
\mathbf{i} & \mathbf{j} & \mathbf{k} \\
1 & 2t & 0 \\
0 & 2 & 0
\end{vmatrix}
= \langle 0, 0, 2 \rangle.
\]

Por lo tanto
\[
||r'(t) \times \vec{r}''(t)|| = 2.
\]

Ahora calculamos la norma de \( r'(t) \)
\[
||r'(t)|| = \sqrt{1 + 4t^2}.
\]

Entonces la curvatura es
\[
\tau(t) = \frac{2}{(1 + 4t^2)^{3/2}}.
\]

Evaluando en $t=1$, obtenemos

Entonces la curvatura es:
\[
\tau(1) = \frac{2}{5^{3/2}}.
\]










\end{document}